% !TEX program = xelatex
% Статья: Сравнение риск-ориентированных методов планирования пути
% четвероногого робота на вероятностной карте проходимости.
% Сборка: tectonic article.tex  (или xelatex article.tex дважды)

\documentclass{misisarticle}

\usepackage{tikz}
\usetikzlibrary{arrows.meta, positioning}

% --- Классификационные коды ---
\grnti{28.23.27}
\asjc{2604, 2207}
\oecd{2.02}

% --- Название ---
\title{Сравнение риск-ориентированных методов планирования пути четвероногого робота на вероятностной карте проходимости}

% --- Авторы ---
\misisauthor{М.~Плешевич}{кафедра инженерной кибернетики НИТУ «МИСиС»}{Москва, Россия}{pleshevich.milena@mail.ru}
\misisauthor{В.~Семёнов}{кафедра инженерной кибернетики НИТУ «МИСиС»}{Москва, Россия}{}

% --- Аннотация и ключевые слова ---
\abstract{%
  Рассматривается задача планирования пути четвероногого робота в недетерминированной среде, представленной вероятностной картой проходимости. Для шагающего робота кратчайший путь не всегда безопасен: уклон поверхности, неровности рельефа и высота препятствий могут привести к потере устойчивости. В работе предложена единая постановка задачи и методика сравнения риск-ориентированных планировщиков: базового алгоритма A* по кратчайшему пути, A* с вероятностно-взвешенной стоимостью, планирования с вероятностным ограничением и планирования по мере риска CVaR. Описан способ перевода уклона, шероховатости и высоты ступени в вероятность прохождения ячейки с учётом ограничений конкретной модели робота. Предложены сценарии на синтетических картах высот с разным уровнем неопределённости и метрики качества: длина маршрута, суммарный риск, время построения пути и доля успешных прохождений в симуляции. Методика позволяет определить, при каком уровне неопределённости карты кратчайший путь перестаёт быть приемлемым, и дать рекомендации по выбору метода и параметра риска.
}

\keywords{%
  четвероногий робот, планирование пути, вероятностная карта проходимости, риск-ориентированное планирование, CVaR, недетерминированная среда, ограничения устойчивости
}

\begin{document}

\maketitle

% ============================================================
\section{Введение}

Четвероногие роботы всё чаще применяются вне лабораторий: для инспекции промышленных объектов, работы в шахтах и на стройплощадках, мониторинга территорий и поисково-спасательных операций. Главное их преимущество перед колёсными платформами~--- способность двигаться по пересечённой местности, ступеням и рыхлому грунту~\cite{miki2022, wellhausen2023}. Однако это преимущество реализуется только тогда, когда система планирования учитывает реальные возможности робота и состояние поверхности.

Классические планировщики пути на сетке, например A*~\cite{hart1968}, ищут маршрут минимальной длины на бинарной карте «свободно/занято». Для шагающего робота такая модель среды слишком груба. Участок с уклоном, близким к предельному, формально проходим, но вероятность падения на нём велика. Кроме того, карта высот, построенная по данным лидара или стереокамеры, содержит ошибки измерения и локализации~\cite{fankhauser2018, thrun2005}, поэтому свойства каждой ячейки известны лишь с некоторой вероятностью. В результате кратчайший путь может проходить по самым опасным участкам.

Естественный способ учесть неопределённость~--- представить среду вероятностной картой проходимости, в которой каждой ячейке сопоставлена вероятность её успешного прохождения, и использовать риск-ориентированный критерий при поиске пути~\cite{fan2021, akella2024}. В литературе предложено несколько семейств таких методов: минимизация ожидаемой стоимости, ограничения на вероятность отказа (chance constraints)~\cite{blackmore2011}, когерентные меры риска, прежде всего условная стоимость под риском CVaR~\cite{rockafellar2000, majumdar2020}, а также графовые представления риска~\cite{lee2025}. Эти методы разработаны для разных платформ и проверены в разных условиях, поэтому из опубликованных результатов трудно понять, какой из них предпочтителен для четвероногого робота и при каком уровне неопределённости карты.

Цель работы~--- сравнить риск-ориентированные методы планирования пути четвероногого робота на вероятностной карте проходимости в единой постановке. Новый алгоритм в статье не предлагается: вклад работы методологический и экспериментальный. Он состоит в следующем:
\begin{itemize}
  \item единая постановка задачи и стенд, в котором методы с разной трактовкой риска реализованы на одной вероятностной карте и в одном симуляторе;
  \item модель перевода уклона, шероховатости и высоты ступени в вероятность прохождения ячейки с учётом допустимого угла наклона и высоты шага конкретной модели робота;
  \item методика анализа зависимости качества планирования от уровня неопределённости карты;
  \item схема получения практических рекомендаций по выбору метода и параметра риска (уровня $\alpha$ для CVaR или порога вероятности) для разных типов местности.
\end{itemize}

Статья организована следующим образом. В разделе~II проведён анализ существующих работ. В разделе~III формально поставлена задача. В разделе~IV описаны построение вероятностной карты, сравниваемые методы, экспериментальный стенд и метрики. Раздел~V содержит выводы и план дальнейших исследований.

% ============================================================
\section{Анализ существующих работ}

\subsection{Оценка проходимости и вероятностные карты}

Стандартным представлением среды для шагающих роботов стала 2,5-мерная карта высот. Библиотека grid\_map~\cite{fankhauser2016} задаёт многослойную регулярную сетку, в слоях которой хранятся высота, её дисперсия и производные признаки. В работе~\cite{fankhauser2018} предложено построение карты высот в системе координат робота с явным учётом неопределённости локализации: для каждой ячейки оценивается не только средняя высота, но и доверительный интервал. Именно такая карта служит входом для вероятностной оценки проходимости.

Проходимость ячейки обычно вычисляют по геометрическим признакам: уклону, шероховатости и высоте перепада~\cite{wermelinger2016}. Пороги этих признаков задаются из возможностей конкретного робота. Альтернативный подход~--- обучать оценку проходимости в симуляции~\cite{chavez2018} или на основе политики локомоции~\cite{frey2022}. Обученные модели точнее описывают реальные возможности робота, но требуют больших объёмов данных и хуже интерпретируются. Для сравнения планировщиков в контролируемых условиях геометрическая модель предпочтительнее: её параметры явно связаны с характеристиками робота, а уровень неопределённости можно задавать непосредственно.

\subsection{Риск-ориентированное планирование}

Наиболее простой способ учесть неопределённость~--- минимизировать ожидаемую стоимость пути. Такой критерий нейтрален к риску: редкие, но тяжёлые исходы (падение робота) усредняются и почти не влияют на решение. Поэтому в робототехнике применяют меры риска, чувствительные к «хвосту» распределения потерь. В работе~\cite{majumdar2020} сформулированы аксиомы, которым должна удовлетворять мера риска для робота, и показано, что им отвечают когерентные меры, в том числе CVaR. Обзор~\cite{akella2024} систематизирует применение мер хвостового риска в планировании, управлении и верификации.

Метод STEP~\cite{fan2021} оценивает проходимость местности как случайную величину и строит путь с минимальным CVaR стоимости. Метод проверен на колёсных и шагающих роботах в рамках проекта DARPA SubT и показал, что учёт хвостового риска снижает число аварийных ситуаций без значительного удлинения пути. В работе~\cite{dixit2021} вместо CVaR использована энтропийная стоимость под риском (EVaR), которая даёт более консервативную оценку и допускает эффективное вычисление.

Альтернативный подход~--- планирование с вероятностным ограничением: путь должен быть кратчайшим при условии, что вероятность отказа не превышает заданного $\delta$~\cite{blackmore2011}. Это интуитивно понятный параметр для пользователя, но точное решение такой задачи на графе требует поиска с ограничением на ресурс, а упрощение до порога на отдельную ячейку теряет гарантии для пути в целом.

Метод TRG-planner~\cite{lee2025} строит разреженный граф риска прохождения: узлы размещаются в областях низкого риска, рёбра взвешиваются интегральным риском. Такой граф ускоряет повторное планирование в неструктурированной среде и был проверен на четвероногом роботе.

\subsection{Планирование для шагающих роботов}

Для шагающих роботов задача планирования обычно разбивается на глобальный уровень (маршрут по карте) и локальный (выбор опорных точек и походки). Система ArtPlanner~\cite{wellhausen2023} сочетает выборочный планировщик с обученной оценкой достижимости и успешно применялась в полевых испытаниях четвероногого робота ANYmal. В~\cite{wermelinger2016} проходимость оценивается по карте высот и используется планировщиком на основе RRT*. Подход R-LGP~\cite{ly2024} учитывает достижимость при совместном планировании задач и движения для мобильного манипулятора. Локальные методы на основе обучения с подкреплением~\cite{kiselev2025} хорошо реагируют на изменения среды, но не дают явной оценки риска глобального маршрута. Качество карты высот, в свою очередь, зависит от точности одометрии и замыкания цикла~\cite{epifanov2025}, и эта погрешность напрямую переходит в неопределённость проходимости.

\subsection{Выводы из анализа}

Существующие работы предлагают отдельные риск-ориентированные методы и проверяют их на собственных платформах, картах и метриках. Систематического сравнения методов с разной трактовкой риска в одинаковых условиях для четвероногого робота найдено не было. Не исследована и зависимость относительного качества методов от уровня неопределённости карты: неясно, при каком шуме имеет смысл отказываться от кратчайшего пути и какой параметр риска выбирать для разных типов местности. Эти вопросы и рассматриваются в настоящей работе.

% ============================================================
\section{Постановка задачи}

Среда задаётся регулярной сеткой размера $W \times H$ с шагом $\Delta$. Для каждой ячейки $i$ известна оценка высоты $\hat z_i$ и её стандартное отклонение $\sigma_i$, характеризующее неопределённость карты. Из карты высот вычисляется вероятность прохождения $p_i \in [0, 1]$~--- вероятность того, что робот пересечёт ячейку без потери устойчивости. Способ её вычисления описан в подразделе~IV-A.

Карта рассматривается как граф $G = (V, E)$, где вершины~--- ячейки, а рёбра соединяют соседние ячейки в 8-связной окрестности. Длина ребра $\ell_{ij}$ равна $\Delta$ или $\sqrt{2}\Delta$. Путь $\pi = (v_0, v_1, \dots, v_n)$ соединяет начальную вершину $v_0 = s$ с целевой $v_n = g$. Его длина
\begin{equation}
  L(\pi) = \sum_{k=1}^{n} \ell_{v_{k-1} v_k}.
  \label{eq:length}
\end{equation}

Будем считать события прохождения разных ячеек независимыми. Тогда вероятность успешного прохождения пути и его риск равны
\begin{equation}
  P(\pi) = \prod_{k=0}^{n} p_{v_k}, \qquad R(\pi) = 1 - P(\pi).
  \label{eq:risk}
\end{equation}
Предположение о независимости упрощает модель: соседние ячейки на самом деле коррелированы, так как ошибка карты имеет пространственную структуру. Влияние этого упрощения оценивается в эксперименте сравнением $R(\pi)$ с наблюдаемой долей неудачных прохождений.

Планировщик $\mathcal{A}$ по карте, вершинам $s$, $g$ и параметру риска $\theta$ строит путь $\pi = \mathcal{A}(M, s, g; \theta)$. Требуется для набора планировщиков $\{\mathcal{A}_1, \dots, \mathcal{A}_K\}$ и набора карт с разным уровнем неопределённости $\sigma_n$ оценить:
\begin{itemize}
  \item длину пути $L(\pi)$;
  \item суммарный риск $R(\pi)$;
  \item время построения пути $T$;
  \item долю успешных прохождений $S$ в симуляции, когда робот движется по пути, а рельеф выбирается из распределения, заданного картой.
\end{itemize}

Исследовательские вопросы работы формулируются так:
\begin{enumerate}[label=\arabic*)]
  \item при каком уровне неопределённости $\sigma_n$ кратчайший путь перестаёт обеспечивать приемлемую долю успешных прохождений;
  \item какой метод даёт лучший компромисс между длиной пути и риском на разных типах местности;
  \item как выбирать параметр риска ($\alpha$ для CVaR, порог вероятности $\delta$ или вес $\lambda$) в зависимости от типа местности и уровня неопределённости.
\end{enumerate}

% ============================================================
\section{Описание методов}

\subsection{Вероятностная карта с учётом устойчивости}

Для каждой ячейки по окрестности радиуса $r$ вычисляются три признака: уклон $\varphi_i$~--- угол между нормалью к аппроксимирующей плоскости и вертикалью; шероховатость $\rho_i$~--- среднеквадратичное отклонение высот от этой плоскости; высота ступени $h_i$~--- максимальный перепад высот между соседними ячейками. Пороги признаков определяются возможностями робота: допустимым углом наклона $\varphi_{\max}$, высотой шага $h_{\max}$ и допустимой шероховатостью $\rho_{\max}$. Для модели Unitree Go2 их значения берутся из технических характеристик и уточняются в симуляции.

Признаки вычисляются по неточной карте и потому сами являются случайными величинами. Будем считать их нормально распределёнными со средними $\hat\varphi_i$, $\hat\rho_i$, $\hat h_i$ и стандартными отклонениями $\sigma^{\varphi}_i$, $\sigma^{\rho}_i$, $\sigma^{h}_i$, полученными распространением неопределённости высот $\sigma_i$. Тогда вероятность того, что признак не превысит порог, равна $\Phi\big((x_{\max} - \hat x_i)/\sigma^{x}_i\big)$, где $\Phi$~--- функция стандартного нормального распределения. Вероятность прохождения ячейки:
\begin{equation}
  p_i = \Phi\!\left(\frac{\varphi_{\max} - \hat\varphi_i}{\sigma^{\varphi}_i}\right)
        \Phi\!\left(\frac{h_{\max} - \hat h_i}{\sigma^{h}_i}\right)
        \Phi\!\left(\frac{\rho_{\max} - \hat\rho_i}{\sigma^{\rho}_i}\right).
  \label{eq:pcell}
\end{equation}
При $\sigma \to 0$ выражение~\eqref{eq:pcell} переходит в бинарную карту проходимости. С ростом неопределённости вероятность ячеек вблизи порога стремится к $0{,}5$, а ячейки, далёкие от порога, остаются надёжно проходимыми или непроходимыми. Ячейки, занятые препятствиями, получают $p_i = 0$. Карта хранится в формате grid\_map~\cite{fankhauser2016} как отдельный слой.

\subsection{Сравниваемые методы}

Все методы реализуются на основе одного поиска A*~\cite{hart1968} на графе $G$ с евклидовой эвристикой и отличаются только стоимостью ребра или правилом отбрасывания вершин. Это исключает влияние деталей реализации на сравнение времени и качества.

\textit{1) Базовый A* по кратчайшему пути.} Карта бинаризуется: ячейка считается свободной при $p_i \ge 0{,}5$. Стоимость ребра равна его длине $\ell_{ij}$. Метод игнорирует риск и служит нижней границей длины пути.

\textit{2) A* с вероятностно-взвешенной стоимостью.} Логарифмируя~\eqref{eq:risk}, получаем $-\ln P(\pi) = \sum_k (-\ln p_{v_k})$. Риск становится аддитивным, и его можно включить в стоимость ребра:
\begin{equation}
  c_{ij} = \ell_{ij} \left(1 - \lambda \frac{\ln p_j}{\Delta}\right),
  \label{eq:wastar}
\end{equation}
где $\lambda \ge 0$~--- вес риска. Нормировка на $\Delta$ делает стоимость независимой от разрешения карты. При $\lambda = 0$ метод совпадает с базовым A*, при больших $\lambda$ ищется путь с максимальной вероятностью успеха.

\textit{3) Планирование с вероятностным ограничением.} Требуется найти кратчайший путь при условии $R(\pi) \le \delta$, то есть $\sum_k (-\ln p_{v_k}) \le -\ln(1-\delta)$. Рассматриваются два варианта. В упрощённом варианте из графа удаляются ячейки с $p_i < p_{\mathrm{th}}$, после чего ищется кратчайший путь; ограничение на путь в целом при этом не гарантируется. В точном варианте используется поиск с ограничением на ресурс: в состояние A* добавляется накопленный риск, а пути, нарушающие ограничение, отбрасываются. Сравнение вариантов показывает, сколько теряется при простом пороге на ячейку.

\textit{4) Планирование по мере риска CVaR.} Следуя~\cite{fan2021}, риск ячейки описывается случайной величиной потерь $\xi_i$, связанной с запасом устойчивости. Для нормального распределения с параметрами $\mu_i$, $s_i$ условная стоимость под риском на уровне $\alpha$ вычисляется в замкнутой форме~\cite{rockafellar2000}:
\begin{equation}
  \mathrm{CVaR}_\alpha(\xi_i) = \mu_i + s_i \frac{\phi\big(\Phi^{-1}(\alpha)\big)}{1 - \alpha},
  \label{eq:cvar}
\end{equation}
где $\phi$~--- плотность стандартного нормального распределения. Стоимость ребра
\begin{equation}
  c_{ij} = \ell_{ij} \big(1 + \lambda\, \mathrm{CVaR}_\alpha(\xi_j)\big).
  \label{eq:cvarcost}
\end{equation}
Параметр $\alpha \in [0, 1)$ задаёт отношение к риску: при $\alpha = 0$ получается ожидаемая стоимость, при $\alpha \to 1$~--- оценка по худшему случаю. Поскольку CVaR~--- когерентная мера риска и обладает субаддитивностью, сумма $\mathrm{CVaR}$ по ячейкам является верхней оценкой CVaR суммарных потерь пути, то есть критерий~\eqref{eq:cvarcost} консервативен.

Основные свойства методов сведены в табл.~\ref{tab:methods}.

\begin{table}[htbp]
  \caption{Сравниваемые методы планирования}
  \label{tab:methods}
  \centering
  \small
  \begin{tabular}{p{0.30\columnwidth}p{0.24\columnwidth}p{0.30\columnwidth}}
    \toprule
    \textbf{Метод} & \textbf{Параметр} & \textbf{Трактовка риска} \\
    \midrule
    A* (кратчайший) & --- & не учитывается \\
    A* с весом риска & $\lambda$ & вероятность отказа пути \\
    Порог на ячейку & $p_{\mathrm{th}}$ & отказ отдельной ячейки \\
    Ограничение на путь & $\delta$ & отказ пути $\le \delta$ \\
    CVaR & $\alpha$, $\lambda$ & хвост распределения потерь \\
    \bottomrule
  \end{tabular}
\end{table}

\subsection{Экспериментальный стенд}

Схема стенда приведена на рис.~\ref{fig:pipeline}. Сравнение проводится на синтетических картах высот размером от $100 \times 100$ до $200 \times 200$ ячеек. Рельеф задаётся шумом Перлина~\cite{perlin2002}, к которому добавляются склоны, ступени и препятствия. Используются четыре типа местности: ровная с препятствиями, холмистая, ступенчатая и смешанная. Для каждой карты истинного рельефа строятся наблюдаемые карты с аддитивным шумом высоты уровня $\sigma_n$, например $\sigma_n \in \{0{,}01;\ 0{,}03;\ 0{,}05;\ 0{,}10\}$~м.

\begin{figure}[htbp]
  \centering
  \begin{tikzpicture}[
      node distance=4mm,
      box/.style={draw, rounded corners=2pt, align=center, font=\footnotesize,
                  minimum width=0.8\columnwidth, minimum height=6mm, inner sep=2pt},
      arr/.style={-{Latex[length=2mm]}}
    ]
    \node[box] (gen)  {Генерация рельефа (шум Перлина,\\ склоны, ступени, препятствия)};
    \node[box, below=of gen]  (noise) {Наблюдаемая карта высот: шум $\sigma_n$};
    \node[box, below=of noise] (prob) {Вероятностная карта проходимости, \eqref{eq:pcell}};
    \node[box, below=of prob]  (plan) {Планировщики: A*, вес риска, ограничение, CVaR};
    \node[box, below=of plan]  (sim)  {Симуляция Unitree Go2 (ROS~2, Gazebo / Isaac Sim)};
    \node[box, below=of sim]   (met)  {Метрики: $L$, $R$, $T$, $S$};
    \draw[arr] (gen) -- (noise);
    \draw[arr] (noise) -- (prob);
    \draw[arr] (prob) -- (plan);
    \draw[arr] (plan) -- (sim);
    \draw[arr] (sim) -- (met);
  \end{tikzpicture}
  \caption{Схема экспериментального стенда для сравнения методов планирования}
  \label{fig:pipeline}
\end{figure}

Симуляция выполняется в ROS~2 с Gazebo или Isaac Sim на модели четвероногого робота Unitree Go2. Робот проходит построенный путь с помощью штатного локального контроллера по истинному рельефу. Прохождение считается неудачным при падении робота, застревании дольше заданного времени или отклонении от пути больше допустимого. Для каждой пары «карта~--- метод» выполняется серия запусков со случайными начальной и целевой точками.

Для проверки на реальном рельефе используется открытый набор данных RELLIS-3D~\cite{jiang2021}. Он записан на колёсной платформе, но содержит облака точек реальной пересечённой местности. По ним строятся карты высот, к которым применяются те же планировщики; оценка ограничивается метриками $L$, $R$ и $T$, так как прохождение моделируется только в симуляции.

\subsection{Метрики и обработка результатов}

Для каждого запуска фиксируются длина пути $L$ по~\eqref{eq:length}, риск $R$ по~\eqref{eq:risk}, вычисленный по истинной карте, время построения $T$ и признак успешного прохождения. Доля успешных прохождений $S$ оценивается по серии запусков с доверительным интервалом Уилсона. Чтобы сравнивать методы на картах разного размера, длина нормируется на длину кратчайшего пути на истинной карте.

Компромисс между длиной и риском анализируется по фронту Парето в координатах $(L, R)$, который строится перебором параметра каждого метода ($\lambda$, $p_{\mathrm{th}}$, $\delta$, $\alpha$). Зависимость $S(\sigma_n)$ для базового A* позволяет найти уровень неопределённости, начиная с которого кратчайший путь становится неприемлемым, например при $S < 0{,}9$. Рекомендации по выбору метода и параметра формируются отдельно для каждого типа местности как точка фронта Парето с минимальной длиной при заданной требуемой доле успешных прохождений.

% ============================================================
\section{Заключение}

В работе поставлена задача сравнения риск-ориентированных методов планирования пути четвероногого робота на вероятностной карте проходимости. Предложена модель, переводящая уклон, шероховатость и высоту ступени с учётом неопределённости карты в вероятность прохождения ячейки по ограничениям конкретного робота. Четыре семейства методов~--- кратчайший путь, взвешивание по вероятности отказа, вероятностное ограничение и CVaR~--- сведены к единой реализации на основе A*, что делает сравнение корректным. Описаны стенд на синтетических картах с управляемым уровнем шума, симуляция на модели Unitree Go2 и метрики качества.

Дальнейшая работа включает реализацию стенда, проведение серии экспериментов, построение фронтов Парето для разных типов местности и проверку на картах, построенных по набору RELLIS-3D. Отдельно планируется оценить влияние предположения о независимости ячеек на точность прогноза риска.

% ============================================================
\begin{thebibliography}{99}

\bibitem{miki2022} Miki T., Lee J., Hwangbo J., Wellhausen L., Koltun V., Hutter M. Learning robust perceptive locomotion for quadrupedal robots in the wild. \textit{Science Robotics}, vol.~7, no.~62, eabk2822, 2022. DOI: 10.1126/scirobotics.abk2822.

\bibitem{wellhausen2023} Wellhausen L., Hutter M. ArtPlanner: Robust legged robot navigation in the field. \textit{Field Robotics}, vol.~3, pp.~413--434, 2023.

\bibitem{hart1968} Hart P.~E., Nilsson N.~J., Raphael B. A formal basis for the heuristic determination of minimum cost paths. \textit{IEEE Transactions on Systems Science and Cybernetics}, vol.~4, no.~2, pp.~100--107, 1968.

\bibitem{fankhauser2018} Fankhauser P., Bloesch M., Hutter M. Probabilistic terrain mapping for mobile robots with uncertain localization. \textit{IEEE Robotics and Automation Letters}, vol.~3, no.~4, pp.~3019--3026, 2018.

\bibitem{thrun2005} Thrun S., Burgard W., Fox D. \textit{Probabilistic Robotics}. MIT Press, Cambridge, 2005.

\bibitem{fan2021} Fan D.~D., Otsu K., Kubo Y., Dixit A., Burdick J., Agha-Mohammadi A. STEP: Stochastic traversability evaluation and planning for risk-aware off-road navigation. \textit{Robotics: Science and Systems (RSS)}, 2021. arXiv:2103.02828.

\bibitem{akella2024} Akella P., Dixit A., Ahmadi M., Lindemann L., Chapman M.~P., Pappas G.~J., Ames A.~D., Burdick J.~W. Risk-aware robotics: Tail risk measures in planning, control, and verification. arXiv:2403.18972, 2024.

\bibitem{blackmore2011} Blackmore L., Ono M., Williams B.~C. Chance-constrained optimal path planning with obstacles. \textit{IEEE Transactions on Robotics}, vol.~27, no.~6, pp.~1080--1094, 2011.

\bibitem{rockafellar2000} Rockafellar R.~T., Uryasev S. Optimization of conditional value-at-risk. \textit{Journal of Risk}, vol.~2, no.~3, pp.~21--41, 2000.

\bibitem{majumdar2020} Majumdar A., Pavone M. How should a robot assess risk? Towards an axiomatic theory of risk in robotics. In: \textit{Robotics Research}, Springer Proceedings in Advanced Robotics, vol.~10, pp.~75--84, 2020.

\bibitem{lee2025} Lee D., Nahrendra I.~M.~A., Oh M., Yu B., Myung H. TRG-planner: Traversal risk graph-based path planning in unstructured environments for safe and efficient navigation. \textit{IEEE Robotics and Automation Letters}, 2025. arXiv:2501.01806.

\bibitem{fankhauser2016} Fankhauser P., Hutter M. A universal grid map library: Implementation and use case for rough terrain navigation. In: \textit{Robot Operating System (ROS): The Complete Reference (Volume~1)}, Springer, pp.~99--120, 2016.

\bibitem{wermelinger2016} Wermelinger M., Fankhauser P., Diethelm R., Krüsi P., Siegwart R., Hutter M. Navigation planning for legged robots in challenging terrain. \textit{2016 IEEE/RSJ International Conference on Intelligent Robots and Systems (IROS)}, pp.~1184--1189, 2016.

\bibitem{chavez2018} Chavez-Garcia R.~O., Guzzi J., Gambardella L.~M., Giusti A. Learning ground traversability from simulations. \textit{IEEE Robotics and Automation Letters}, vol.~3, no.~3, pp.~1695--1702, 2018.

\bibitem{frey2022} Frey J., Hoeller D., Khattak S., Hutter M. Locomotion policy guided traversability learning using volumetric representations of complex environments. \textit{2022 IEEE/RSJ International Conference on Intelligent Robots and Systems (IROS)}, 2022.

\bibitem{dixit2021} Dixit A., Ahmadi M., Burdick J.~W. Risk-sensitive motion planning using entropic value-at-risk. \textit{2021 European Control Conference (ECC)}, 2021.

\bibitem{ly2024} Ly K.~T., Semenov V., Risiglione M., Merkt W., Havoutis I. R-LGP: A reachability-guided logic-geometric programming framework for optimal task and motion planning on mobile manipulators. \textit{2024 IEEE International Conference on Robotics and Automation (ICRA)}, Yokohama, pp.~14917--14923, 2024. DOI: 10.1109/ICRA57147.2024.10611389.

\bibitem{kiselev2025} Киселев С.~Д., Алтунян А.~В. Локальные методы планирования траекторий на основе обучения с подкреплением. \textit{Искусственный интеллект в промышленных, коммерческих, медицинских и финансовых приложениях, 2025}: сб. ст. науч.-техн. семинара. Вып.~3 / под ред. А.~Р.~Ефимова. М.: НИТУ «МИСИС», 2025. С.~84--89.

\bibitem{epifanov2025} Епифанов К.~А. Исследование алгоритмов замыкания цикла в лидарной одометрии. \textit{Искусственный интеллект в промышленных, коммерческих, медицинских и финансовых приложениях, 2025}: сб. ст. науч.-техн. семинара. Вып.~3 / под ред. А.~Р.~Ефимова. М.: НИТУ «МИСИС», 2025. С.~57--61.

\bibitem{perlin2002} Perlin K. Improving noise. \textit{ACM Transactions on Graphics}, vol.~21, no.~3, pp.~681--682, 2002.

\bibitem{jiang2021} Jiang P., Osteen P., Wigness M., Saripalli S. RELLIS-3D dataset: Data, benchmarks and analysis. \textit{2021 IEEE International Conference on Robotics and Automation (ICRA)}, pp.~1110--1116, 2021.

\end{thebibliography}

\end{document}
